We have also compared the simulated sea ice with satellite observations (Figure~\ref{fR:seaice}). The observations are the OSI SAF sea ice concentration climate data record (version 3.0) averaged over 1979–2008, the first 30 years of the satellite era. The model fields are the last 100 years of the pre-industrial run.

\rfig{0.9}{../letter_figs/figR21_seaice_obs.png}{New analysis. Southern Ocean sea ice concentration in the pre-industrial simulation (last 100 years) and in the OSI SAF satellite record for 1979–2008 (sea ice concentration climate data record, version 3.0). (a, b) September, the month with the most ice, with the 15\% ice edge of the model (orange, dashed) and of the observations (black). (c) Model minus observations in September. (d, e) February, the month with the least ice. (f) Seasonal cycle of the Southern Hemisphere sea ice extent (area with a concentration of at least 15\%, solid lines) and sea ice area (dashed lines), with the range of individual years shaded.\label{fR:seaice}}

The model and the observations agree in the timing of the seasonal cycle, with the most ice in September and the least in February. They also agree on the position of the winter ice edge in the Ross Sea and Amundsen–Bellingshausen sector (180°–60°W), where the simulated September edge lies within 0.4° of latitude of the observed one. Elsewhere the model has too little ice. Its September extent is 9.9 million km\textsuperscript{2}, about half of the observed 19.1 million km\textsuperscript{2}, and its February extent is 0.7 against 3.3 million km\textsuperscript{2}, so almost all of the ice melts in summer. Most of the winter deficit is in the Weddell sector. There the simulated September ice edge lies about 10° of latitude south of the observed one, and between 60 and 70°S and 30°W and 30°E the simulated September ice concentration is 8\% against 96\% in the observations. This ice-free part of the Weddell gyre is where the model convects (Figure~\ref{fR:mld}). Too little ice also leaves a larger area of open water for the atmosphere to cool, which is consistent with the strong thermal contribution in our pre-industrial transformation.
